\section{Implementation and numerical verification}
\label{sec:implementation}

\subsection{Python architecture}

The entry point \file{solve} accepts a finite 2D observation, a centered unit-sum PSF, and a method identifier. It returns a floating-point estimate on the same grid. The input contract catches shape mismatch, unsupported algorithm names, nonfinite values, invalid parameter ranges, negative PSF coefficients, and negative observations in the multiplicative solvers. Conversion to a display range happens at scoring or visualization. It is separate from the reconstruction itself.

\begin{table}[!htbp]
\centering
\caption{Implementation components and their normal entry points relative to the project root.}
\small
\begin{tabularx}{\linewidth}{@{}p{.38\linewidth}Y@{}}
\toprule Component & Responsibility\\\midrule
\file{code/deconv/core.py} & TIFF/PNG handoff, PSF construction, circular blur, and measurement simulation.\\
\file{code/deconv/solvers.py} & Direct spectral filters, iterative estimators, and the RL rule used for matching the plugin.\\
\file{code/deconv/metrics.py} & MSE, fixed-peak PSNR, and weighted local SSIM.\\
\file{code/scripts/} & Dataset preparation, experiment orchestration, audit, and report inputs.\\
\file{code/imagej/BatchReference.java} & ImageJ TIFF reading, plugin execution, independent metric calculation, and CSV export.\\
\file{code/tests/test_deconv.py} & Analytical examples and independent numerical oracles.\\
\bottomrule
\end{tabularx}
\end{table}

Listing~\ref{lst:forward-code} shows how the forward operator places the PSF origin at index \((0,0)\) before computing the OTF. The inverse \file{irfft2} receives the original image shape so it also reconstructs an odd last dimension correctly. The solver uses the same OTF for blurring and its complex conjugate for back-projection. Listing~\ref{lst:spectral-code} computes the Laplacian spectrum directly from the row and column frequencies, including for one-row and one-column inputs.

\reportlisting{Python}{deconv/core.py}{46}{63}{
  caption={PSF centering, validation, and circular convolution in \texttt{core.py}.},
  label={lst:forward-code}
}

\subsection{Iterative update and stability checks}

Listing~\ref{lst:iterative-code} shows the RL and RLTV loop. The main branch clamps negative roundoff in the prediction and correction and uses \(\varepsilon=10^{-12}\). The plain RL matching branch uses the plugin's observation initialization and zero-ratio threshold. RLTV checks the denominator before division and verifies that every updated value is finite. Appendix~\ref{app:code} includes the full \file{solve} implementation and the experiment and scoring excerpts.

For \(Q=MN\), the direct methods take \(O(Q\log Q)\) time and \(O(Q)\) working memory. RL and RLTV take \(O(KQ\log Q)\) time, with \(O(Q)\) extra work per iteration for the TV terms. The constant can be substantial because the prediction and correction each require a forward and inverse transform. Reusing the PSF spectrum reduces unnecessary FFT calls without changing the iteration count. The code distinguishes the Python float64 arithmetic from the float32 handoff to ImageJ.

\subsection{ImageJ reference workflow}

The Java batch launcher invokes \file{DeconvolutionLab2_Run} from DeconvolutionLab2 2.1.2 on ImageJ 1.54p \citep{sage,dl2docs}. Listing~\ref{lst:imagej-code} shows the plugin invocation and scoring handoff. Appendix~\ref{app:imagej} gives the corresponding settings for replay in the GUI.

The launcher sets macro options rather than typing into a GUI. It reads exported TIFFs through ImageJ, checks their bit depth and dimensions, and independently computes the metrics. The reference runtime is Java \JVMVersion. The batch records the command for every case, the JAR URLs and SHA-256 hashes, and the algorithm outputs as float TIFF. It also calls the plugin's \file{Assessment.psnr} for an additional native-score check. The JAR and the Python source may use different single-precision normalization and discrete derivative rules, so the settings and source comparison are part of the evidence \citep{dl2code}.

\subsection{Verification strategy}

The 20 Python tests cover hand calculations, agreement with spatial convolution for an asymmetric PSF, odd/even centering, a dense Tikhonov oracle, the adjoint in RL, a directional derivative of TV, SSIM from an explicit window, reproducible simulation, raw TIFF round-trip, and equivalence of the two FFT branches. The one-row Laplacian test protects the direct spectral implementation against support-size errors. A separate one-step oracle tests the plugin's zero-ratio convention.

The main test audit checks the exact key set for 240 Python records, then reads every referenced TIFF and recomputes the metrics. The maximum metric difference due to the float32 export is \PythonMetricError. A second batch verifies the 240 main plugin metric records plus two for an asymmetric impulse check. The impulse is placed at the border so a wrong phase or missing wrap is observable. Its largest difference from the Python result is \ProbeError. For plugin scoring, both implementations read the same saved pixels. This makes the metric comparison logically independent of the restoration comparison.

Computational measurements time only \file{solve}, excluding TIFF handling, metric calculation, figures, and JVM overhead. They alternate the two FFT modes after warm-up on one validation input. A separate \file{tracemalloc} pass records tracked allocations. This design supports statements about those Python implementations on that machine, while Java serves as a numerical reference.
